\section{Un look-ahead soft a due passi}
\label{sec:softq}

Questa sezione propone una variante di BSP.
È la sola parte algoritmica del lavoro.
La variante usa le equazioni di Bellman soft del Soft Actor-Critic (SAC) di Haarnoja e coautori~\citep{haarnoja2018soft} e le combina con l'identità locale della Proposizione~\ref{prop:local}.
Il ponte con il resto del lavoro è la somma termodinamica di Shannon: SAC la applica alle azioni della ricerca, come SP la applica ai warning e il potenziale 1RSB agli stati.
Gli enunciati algebrici di questa sezione sono verificati in Lean~4 con Mathlib, nel file \texttt{lean/SoftQ.lean} del codice.
Fra parentesi quadre diamo il nome del teorema Lean corrispondente.

\subsection{Il problema di controllo}

Lo stato è la formula residua $F$, con il suo punto fisso di SP.
Un'azione è una coppia $(k,s)$: si fissa la variabile libera $k$ al valore $s$.
La transizione è deterministica e porta nella formula ridotta $F^{(k,s)}=F|_{x_k=s}$, dopo la propagazione unitaria.
La ricompensa $r(F,(k,s))$ sarà definita nella~\eqref{eq:tsallis-reward}.
Una policy $\pi(\cdot\mid F)$ è una distribuzione sulle azioni.
SAC massimizza la somma delle ricompense più l'entropia della policy, pesata con una temperatura $\alpha>0$.
Con transizioni deterministiche le equazioni di Bellman soft sono
\begin{equation}
Q_\alpha(F,a)=r(F,a)+V_\alpha(F^{a}),\qquad
V_\alpha(F)=\alpha\log\sum_a e^{Q_\alpha(F,a)/\alpha},\qquad
\pi_\alpha(a\mid F)=e^{[Q_\alpha(F,a)-V_\alpha(F)]/\alpha}.
\label{eq:soft-bellman}
\end{equation}
La policy di SAC è quindi la distribuzione di Gibbs dei valori $Q_\alpha$, e vale l'identità $Q_\alpha-V_\alpha=\alpha\log\pi_\alpha$.

\begin{proposizione}[Il valore soft è una somma termodinamica]
\label{prop:sac-shannon}
Sia $Q_a$ un numero reale per ogni azione $a$ di un insieme finito, e sia $\alpha>0$.
Per ogni distribuzione $p$ sulle azioni vale
\begin{equation}
\sum_ap_aQ_a+\alpha H(p)\le\alpha\log\sum_ae^{Q_a/\alpha},
\label{eq:sac-bound}
\end{equation}
e l'uguaglianza vale per $p_a\propto e^{Q_a/\alpha}$.
[\texttt{sac\_objective\_le\_soft}, \texttt{sac\_objective\_gibbs}]
\end{proposizione}

\begin{proof}
È la Proposizione~\ref{prop:shannon} con $x_a=-Q_a$ e $T=\alpha$, letta con il segno opposto.
\end{proof}

Il valore $V_\alpha$ è quindi la somma di Maslov $\oplus_\alpha$ della Sezione~\ref{sec:notation} applicata ai valori delle azioni, cioè $V_\alpha=-\bigoplus_{\Sh,\alpha}\{-Q_a\}_a$.
Per $\alpha\to0$ esso tende a $\max_aQ_a$, e la policy diventa la scelta deterministica della migliore azione.
Nel lavoro la somma di Shannon compare quindi su tre livelli.
Il primo è l'aggregazione dei warning (Sezione~\ref{sec:spy}).
Il secondo è il potenziale dei cluster, dove la temperatura è $1/y$.
Il terzo è la scelta delle mosse della ricerca, dove la temperatura è $\alpha$.
Per l'Osservazione~\ref{oss:no-free-T} i primi due livelli hanno la stessa temperatura.
Il terzo livello non ha questo vincolo: $\alpha$ è la temperatura di azione $\Tact$ della Sezione~\ref{sec:response}, e non corrisponde a nessun potenziale 1RSB.

\subsection{Tre ostacoli a una traduzione letterale}

\begin{proposizione}[La complessità telescopa]
\label{prop:telescoping}
Sia $F_0,F_1,\dots,F_n$ una traiettoria della decimazione e sia $r_t=\Sig(F_{t+1})-\Sig(F_t)$.
Allora $\sum_{t<n}r_t=\Sig(F_n)-\Sig(F_0)$.
Due traiettorie con gli stessi estremi hanno la stessa somma.
[\texttt{telescoping}, \texttt{same\_return}]
\end{proposizione}

Una ricompensa di questa forma è una differenza di potenziali.
Ng, Harada e Russell~\citep{ng1999policy} mostrano che una ricompensa di questo tipo non cambia la policy ottima.
La richiesta di una discesa lenta della complessità è quindi una richiesta sulla forma della traiettoria, e non sulla perdita totale.

\begin{osservazione}[Il look-ahead a un passo]
\label{oss:one-step}
Si fissi $x_k=s$ e si riconverga SP.
Per le Proposizioni~\ref{prop:local} e~\ref{prop:remainder} la variazione misurata della complessità è $\log(1-w_k^{-s})+R_{k,s}$, con $R_{k,s}$ di secondo ordine.
Un look-ahead a un passo sulla complessità ricalcola quindi il bias $b_k$, a meno di un termine di secondo ordine, al prezzo di una riconvergenza completa.
\end{osservazione}

Il secondo passo letterale di SAC è $V_\alpha(F^{(k,s)})=\alpha\log\sum_{l,t}e^{r_l^t/\alpha}$, dove la somma corre sulle mosse $(l,t)$ disponibili dopo la prima.
La proposizione seguente mostra che anche questo valore porta poca informazione.

\begin{proposizione}[Il secondo passo letterale degenera]
\label{prop:degenerate}
\begin{enumerate}
\item Siano $b_l\in[0,1]$ per ogni $l$ di un insieme finito, e sia $b_{l_0}=1$ per un $l_0$. Allora $\max_l\log b_l=0$. [\texttt{hard\_two\_step\_degenerate}]
\item Siano $Z>0$, $Z+\delta>0$ e $\alpha\ge0$. Allora
\begin{equation}
\frac{\alpha\delta}{Z+\delta}\le\alpha\log(Z+\delta)-\alpha\log Z\le\frac{\alpha\delta}{Z}.
\label{eq:dilution}
\end{equation}
[\texttt{dilution}]
\end{enumerate}
\end{proposizione}

\begin{proof}
Il primo punto segue da $\log b_l\le0$ e $\log b_{l_0}=0$.
Il secondo segue da $1-1/x\le\log x\le x-1$ con $x=(Z+\delta)/Z$.
\end{proof}

Nel limite $\alpha\to0$ il secondo passo vale $\max_{l}\log b_l$.
Durante quasi tutta la decimazione esiste una variabile con $b_l=1$, congelata o libera in tutti i cluster, e il primo punto dà zero per ogni prima mossa.
Per $\alpha>0$ la somma $Z=\sum_{l,t}e^{r_l^t/\alpha}$ ha $2N_t$ termini.
La prima mossa cambia solo i termini delle variabili vicine a $k$, quindi $\delta=O(1)$.
Se una frazione finita dei termini è di ordine uno, allora $Z=\Theta(N_t)$, e la~\eqref{eq:dilution} dà una differenza di ordine $1/N_t$ fra due prime mosse.

\subsection{La ricompensa di Tsallis}

Sia $\lnq x=(x^{1-q}-1)/(1-q)$ il logaritmo di Tsallis della Proposizione~\ref{prop:tsallis}.
Definiamo la ricompensa della mossa $(k,s)$ come
\begin{equation}
r_k^s=\lnq\bigl(1-w_k^{-s}\bigr).
\label{eq:tsallis-reward}
\end{equation}
Per $q=1$ essa è la variazione della complessità data dall'identità locale~\eqref{eq:local-identity}.

\begin{proposizione}[Proprietà della ricompensa]
\label{prop:tsallis-reward}
Siano $x,y>0$ e $q\ne1$.
\begin{enumerate}
\item $\lnq(xy)=\lnq x+\lnq y+(1-q)\lnq x\,\lnq y$. [\texttt{lnq\_mul}]
\item Se $q>1$ e $x,y<1$, allora $\lnq(xy)<\lnq x+\lnq y$. [\texttt{lnq\_split}]
\item Se $q>1$, la funzione $\lnq$ è strettamente crescente. [\texttt{lnq\_strictMono}]
\item Per $q\to1$ si ha $\lnq x\to\log x$. [\texttt{lnq\_tendsto\_log}]
\end{enumerate}
\end{proposizione}

\begin{proof}
Il primo punto segue da $(xy)^{1-q}=x^{1-q}y^{1-q}$.
Per il secondo, con $q>1$ e $x<1$ si ha $x^{1-q}>1$ e quindi $\lnq x<0$; lo stesso vale per $y$, e il termine $(1-q)\lnq x\,\lnq y$ è negativo.
Il terzo segue da $x^{1-q}$ decrescente e da $1-q<0$.
Il quarto è la derivata di $t\mapsto x^t$ in $t=0$, con $t=1-q$.
\end{proof}

Il primo punto dice che per $q\ne1$ la somma delle ricompense non telescopa.
Il secondo dice che per $q>1$ una caduta di un fattore $xy$ in un solo passo costa più di due cadute di fattori $x$ e $y$.
La ricompensa premia quindi una discesa lenta e regolare.
Il terzo dice che, su un solo passo, la ricompensa ordina le mosse come il bias $b_k$.
La deformazione di Tsallis agisce qui sulla ricompensa della ricerca.
Nella Sezione~\ref{sec:qaxis} essa agiva invece sui pesi dei cluster, e le due deformazioni non vanno confuse.

\subsection{Il valore a campi congelati}

Dopo la prima mossa approssimiamo il futuro con i campi congelati.
Questo significa che la ricompensa di una mossa futura $(l,t)$ è $r_l^t$ calcolata con le grandezze $w_l^\pm$ di $F^{(k,s)}$, e non dipende dalle mosse fatte prima di essa.

\begin{proposizione}[Fattorizzazione]
\label{prop:frozen}
Sia $L$ un insieme finito di variabili, e siano $r_l^\pm$ numeri reali.
\begin{enumerate}
\item Il guadagno di un assegnamento completo $t\in\{\pm1\}^L$ è $\sum_lr_l^{t_l}$, per ogni ordine in cui le variabili vengono fissate. [\texttt{frozen\_order\_free}]
\item Per ogni $\alpha>0$ vale
\begin{equation}
\alpha\log\sum_{t\in\{\pm1\}^L}\exp\Bigl(\frac1\alpha\sum_{l\in L}r_l^{t_l}\Bigr)=\sum_{l\in L}v_\alpha(l),\qquad
v_\alpha(l)=\alpha\log\bigl(e^{r_l^+/\alpha}+e^{r_l^-/\alpha}\bigr).
\label{eq:frozen-value}
\end{equation}
[\texttt{frozen\_soft\_value}]
\end{enumerate}
\end{proposizione}

\begin{proof}
Il primo punto è l'invarianza di una somma finita per permutazione.
Il secondo segue da $\sum_t\prod_le^{r_l^{t_l}/\alpha}=\prod_l\sum_{t_l}e^{r_l^{t_l}/\alpha}$ e dal logaritmo di un prodotto.
\end{proof}

Il valore del futuro a campi congelati è quindi
\begin{equation}
V_\alpha(F)=\sum_{l\ \text{libera in}\ F}v_\alpha(l).
\label{eq:frozen-V}
\end{equation}
Abbiamo tolto il fattore $N_t!$ degli ordinamenti, perché ordini diversi portano allo stesso assegnamento.
Il valore~\eqref{eq:frozen-V} è estensivo, a differenza del secondo passo letterale.

\begin{proposizione}[Limite duro del valore di sito]
\label{prop:site-limit}
Per $\alpha>0$ vale $\max(r_l^+,r_l^-)\le v_\alpha(l)\le\max(r_l^+,r_l^-)+\alpha\log2$.
Quindi $v_\alpha(l)\to\max(r_l^+,r_l^-)=\lnq b_l$ per $\alpha\to0^+$, se $q\ge1$.
[\texttt{soft2\_ge\_max}, \texttt{soft2\_le\_max\_add}, \texttt{soft2\_tendsto\_max}]
\end{proposizione}

\begin{proof}
La somma $e^{r_l^+/\alpha}+e^{r_l^-/\alpha}$ sta fra $e^{m/\alpha}$ e $2e^{m/\alpha}$, con $m=\max(r_l^+,r_l^-)$.
L'uguaglianza $m=\lnq b_l$ segue dalla monotonia di $\lnq$ e da $b_l=1-\min(w_l^+,w_l^-)$.
\end{proof}

\subsection{Il punteggio}

Il valore soft della mossa $(k,s)$ è $Q_\alpha(F,(k,s))=r_k^s+V_\alpha(F^{(k,s)})$.
Siano $L$ e $L'$ gli insiemi delle variabili libere in $F$ e in $F^{(k,s)}$.
L'insieme $L'$ non contiene $k$ e le variabili fissate dalla propagazione unitaria.

\begin{proposizione}[Decomposizione del punteggio]
\label{prop:score}
Vale
\begin{equation}
Q_\alpha(F,(k,s))-V_\alpha(F)=\alpha\log\pi_\alpha(s\mid k)+D_\alpha(k,s),
\label{eq:score}
\end{equation}
con
\begin{equation}
\pi_\alpha(s\mid k)=\frac{e^{r_k^s/\alpha}}{e^{r_k^+/\alpha}+e^{r_k^-/\alpha}},\qquad
D_\alpha(k,s)=\sum_{l\in L'}v_\alpha\bigl(l;F^{(k,s)}\bigr)-\sum_{l\in L\setminus k}v_\alpha(l;F).
\label{eq:D}
\end{equation}
[\texttt{score\_decomposition}, \texttt{sac\_identity}]
\end{proposizione}

\begin{proof}
Si separa il termine di $k$ in $V_\alpha(F)$, e si usa $\alpha\log\pi_\alpha(s\mid k)=r_k^s-v_\alpha(k)$.
\end{proof}

Il primo termine è l'identità di SAC applicata alla scelta del segno di $k$.
Il secondo termine è il vero look-ahead: misura quanto la mossa cambia il valore soft delle altre variabili.

\begin{corollario}[Variabili congelate e variabili libere]
\label{cor:frozen-free}
Se $k$ non è congelata in nessun cluster, cioè $w_k^+=w_k^-=0$, allora $v_\alpha(k)=\alpha\log2$ e $\alpha\log\pi_\alpha(s\mid k)=-\alpha\log2$.
Se $w_k^{-s}=0$, allora $r_k^s=0$ e $0\le v_\alpha(k)\le\alpha e^{r_k^{-s}/\alpha}$.
Se inoltre $w_k^s\to1$, cioè $k$ è congelata al valore $s$ in quasi tutti i cluster, allora $r_k^{-s}\to-\infty$ e $\alpha\log\pi_\alpha(s\mid k)\to0$.
[\texttt{soft2\_unfrozen}, \texttt{soft2\_frozen}]
\end{corollario}

Le due variabili del corollario hanno entrambe $b_k=1$, e la Sezione~\ref{sec:response} ha mostrato che il criterio della complessità non le distingue.
Il punteggio soft le distingue: per $\alpha>0$ esso fissa prima le variabili congelate e lascia libere le variabili libere in tutti i cluster.
Questo è lo stesso ordine del secondo criterio di BSP, la polarizzazione $|w_k^+-w_k^-|$.

\begin{osservazione}[Il termine a un passo si cancella]
\label{oss:cancel}
Per $\alpha\to0$ e $s$ il valore più probabile si ha $\alpha\log\pi_\alpha(s\mid k)\to0$, e la~\eqref{eq:score} si riduce a $D_0(k,s)$.
Con campi congelati il costo $\lnq b_k$ si paga comunque, prima o dopo, e il punteggio dipende solo dall'effetto della mossa sulle altre variabili.
Per questo motivo teniamo il bias di BSP come filtro dei candidati.
\end{osservazione}

La somma in~\eqref{eq:D} corre su tutte le variabili libere, ma riceve contributi soprattutto dalle variabili vicine a $k$.
Il Corollario~\ref{cor:tree} e la Proposizione~\ref{prop:remainder} suggeriscono questa località, ma essa è un'ipotesi da misurare.
Se la località vale, $D_\alpha$ è di ordine uno e non si diluisce come il secondo passo letterale.

\subsection{L'algoritmo SQ$(M,\alpha,q)$}

La variante cambia solo la scelta delle variabili da fissare in un passo di decimazione.
Sia $B=\max(1,\lfloor fN_t\rfloor)$ la dimensione del batch di BSP.
\begin{enumerate}
\item Si prendono come candidati le prime $M$ variabili libere nell'ordine di BSP. Se $M\le B$ l'algoritmo è BSP.
\item Per ogni candidato $k$ con il valore più probabile $s$ si esegue una sonda: si fissa $x_k=s$, si applica la propagazione unitaria e si riconverge SP su una copia della formula. Se la sonda trova una contraddizione o non converge, il punteggio vale $-\infty$. Altrimenti si calcola $V_\alpha(F^{(k,s)})$.
\item Il punteggio è il membro destro della~\eqref{eq:score}.
\item Per $\alpha=0$ si fissano i $B$ candidati con il punteggio più alto. Per $\alpha>0$ si estraggono $B$ candidati senza reinserimento con probabilità proporzionali a $e^{\text{punteggio}/\alpha}$, con il metodo Gumbel-top-$k$~\citep{kool2019stochastic}.
\item I passi di backtracking restano quelli di BSP.
\end{enumerate}
Per $M\le B$ si ritrova BSP.
Nel codice abbiamo verificato che questo limite, e il caso in cui le sonde vengono eseguite ma l'ordine di BSP viene mantenuto, riproducono BSP passo per passo.
Il costo è di $M$ riconvergenze di SP per ogni passo di decimazione.
Il look-ahead a un passo dell'Osservazione~\ref{oss:one-step} è il caso in cui il punteggio è la complessità dopo la sonda.

SAC stima $Q_\alpha$ con una rete e con l'apprendimento.
Qui il modello della dinamica è noto, perché è SP stessa, e $Q_\alpha$ si calcola con una sonda.
La famiglia di algoritmi più vicina è quindi la pianificazione a massima entropia~\citep{xiao2019maximum}, più che l'actor-critic.
